\documentclass[11pt]{article}
\usepackage[T1]{fontenc}
\usepackage{lmodern}
\usepackage{amsmath,amssymb,amsthm}
\usepackage[margin=1in]{geometry}
\usepackage[colorlinks=true,linkcolor=blue,citecolor=blue,urlcolor=blue]{hyperref}
\newtheorem{theorem}{Theorem}
\newtheorem{lemma}[theorem]{Lemma}
\newcommand{\Prob}{\mathbb P}
\newcommand{\E}{\mathbb E}
\newcommand{\capT}{\operatorname{cap}_T}
\title{The exponential excursion-height rate of planar Ising disks}
\author{Kaiwen Jin\\\small\texttt{yui.ui.siki@gmail.com}}
\date{October 13, 2026}

\begin{document}
\maketitle
\begin{abstract}
We determine the exponential rate at which a topology-preserving planar
Ising disk reaches a large area during one excursion from a single square.
For addition rate $\beta$, deletion rate $1$, and every
$0<\beta<1/\kappa$, the excursion maximum $\xi$ satisfies
$\lim_{T\to\infty}\Prob(\xi\ge T)^{1/T}=\kappa\beta$, where $\kappa$ is
the area growth constant of square-lattice self-avoiding polygons.
This proves Conjecture 7.11 of Baryshnikov and Onaran.
The lower bound comes from concatenating arbitrary fixed-area disks with
one connector square and distributing a unit flow over the resulting
growth paths. Cumulative column areas recover each completed prefix,
which bounds congestion without requiring disjoint paths within a block.
\end{abstract}

\section{The model and the result}

The topology-preserving dynamics of Baryshnikov and Onaran~\cite{BO}
samples finite square-lattice disks by adding and deleting squares while
preserving both contractibility and local regularity. Their planar
ergodicity theorem identifies the subcritical fugacity range. For rare
large excursions, Lemma 7.10 of~\cite{BO} gives an upper bound on the
exponential tail rate, and Conjecture 7.11 asks whether this bound is exact.
The question concerns an excursion probability rather than an equilibrium
count: many shapes may be reached through the same small set of transitions.

For $z\in\mathbb Z^2$, let $Q_z=z+[-1/2,1/2]^2$. A finite collection
$C$ of these squares is a \emph{clump} if its geometric union $|C|$ is
contractible and regular. Here regularity means that every $x\in|C|$ is
either an interior point or has
$(B_\varepsilon(x)\cap|C|)\setminus\{x\}$ contractible for all sufficiently
small $\varepsilon>0$. Write $\#C$ for its number of squares and
$o=\{Q_0\}$ for the initial state.

A deletion of $Q\ne Q_0$ is legal when $C\setminus\{Q\}$ is a clump and
there is a strong deformation retraction $H$ of $|C|$ onto
$|C\setminus\{Q\}|$ with $H(Q\times[0,1])\subseteq Q$. Addition is the
inverse move. Starting at $o$, let $C_t$ be the continuous-time chain with
rate $\beta$ for each legal addition and rate $1$ for each legal deletion.
Let $\zeta$ be its first return to $o$ after departure, and define
\[
 \xi=\sup_{0\le t<\zeta}\#C_t,\qquad p_T=\Prob(\xi\ge T).
\]
Let $a_n$ count square-lattice self-avoiding polygons of area $n$ up to
translation, with orientations retained, and let
$\kappa=\lim_{n\to\infty}a_n^{1/n}$ be their area growth constant.

\begin{theorem}\label{thm:main}
For every fixed $0<\beta<1/\kappa$,
\begin{equation}\label{eq:rate}
 \lim_{\substack{T\to\infty\\T\in\mathbb N}}p_T^{1/T}=\kappa\beta.
\end{equation}
Moreover, with $\mu=\E\zeta\in(0,\infty)$ and $\eta_T=\mu/p_T$,
$\lim_{T\to\infty,\,T\in\mathbb N}\eta_T^{1/T}=1/(\kappa\beta)$.
\end{theorem}

Thus the area entropy and the fugacity give the exact exponential scale
for large excursions. The assertion is pointwise in $\beta$ throughout
the open subcritical interval; it does not prescribe a prefactor or a
uniform rate of convergence in $\beta$.

We use two planar facts from~\cite{BO}. Theorem 5.1 states that every
nonsingleton clump has two distinct squares whose deletions leave clumps
through strong deformation retractions supported inside the removed
squares. This geometric statement imposes no distinguished-square
exclusion. The source homotopy condition holds for $t<1$; closedness and
continuity give it at $t=1$ as well. A clump containing $Q_0$ can therefore
be reduced by legal deletions to $o$: at each stage one of the two
geometrically deletable squares differs from $Q_0$. Reversing this sequence
gives a legal monotone growth path retaining $Q_0$. Remark 6.3 and equation (4)
of~\cite{BO} identify unrooted planar clumps with the polygons counted by
$a_n$ and bound the number of origin-containing size-$n$ clumps by $na_n$.

Theorem 6.5 of~\cite{BO} gives nonexplosion and ergodicity in the stated
range, and its Definition 7.2 records the resulting positive finite mean
return time. The source parameters $p_\downarrow=1$ and
$p_\uparrow=\beta$ give exactly our rates. They are admissible since a
straight strip gives $a_n\ge1$, hence $\kappa\ge1$ and $\beta<1$.
We use the defining growth limit for $\kappa$; no additional enumeration
inequality is needed.

\section{Excursions and conductance}

Fix an integer $T\ge2$. Retain all origin-containing clumps with at most
$T$ squares and all legal edges between them. Regard the size-$T$ states
as the terminal set $A_T$. This is a finite graph, and every nonterminal
state has a path to $o$ by the deletion argument above. Give an edge from
a size-$n$ state to a size-$(n+1)$ state the conductance
\begin{equation}\label{eq:conductance}
 c(C,C')=\beta^n.
\end{equation}
Indeed, the reversible weight $w(C)=\beta^{\#C-1}$ gives this conductance
from either direction using the specified rates. At states below size $T$,
the conductance ratios are exactly the actual jump probabilities of the
original chain stopped at $o$ or $A_T$.

Let $v$ be $1$ at $o$, $0$ on $A_T$, and, at other states, the probability
of hitting $o$ before $A_T$. Absorption is almost sure: in the finite
stopped graph, each state has a positive-probability path to an absorbing
state, and a uniform positive lower bound over these finitely many paths
can be applied in repeated bounded blocks of jumps. The first-step
equation makes $v$ harmonic off the boundary. Put
\[
 \capT=\sum_{C\sim o}c(o,C)(1-v(C)).
\]
There are exactly four legal exits from $o$, its four edge-neighbor
additions, each of conductance $\beta$. Corner-only contacts are not
regular. Conditioning on the first actual jump therefore gives
\begin{equation}\label{eq:escape}
 p_T=\frac{\capT}{4\beta}.
\end{equation}
Holding times do not affect this event. Finite summation by parts gives
\begin{equation}\label{eq:voltage-energy}
 \capT=\sum_{\{x,y\}}c(x,y)(v(x)-v(y))^2,
\end{equation}
where undirected edges are counted once: the interior terms vanish by
harmonicity, the terminal terms by $v=0$, and only the term at $o$ remains.

For completeness, the two variational bounds needed below follow directly
from this identity. If $f(o)=1$ and $f=0$ on $A_T$, writing $f=v+g$ and
summing the cross term by parts gives
\[
 \sum_{\{x,y\}}c(x,y)(f(x)-f(y))^2
 =\capT+\sum_{\{x,y\}}c(x,y)(g(x)-g(y))^2\ge\capT.
\]
Choose $f=1$ below size $T$ and $f=0$ at size $T$. There are at most
$T$ downward edges from each of at most $Ta_T$ terminal states, so
\begin{equation}\label{eq:upper}
 \capT\le T^2a_T\beta^{T-1},\qquad
 p_T\le\tfrac14 T^2a_T\beta^{T-2}.
\end{equation}
Taking $T$th roots proves
$\limsup_{T\to\infty}p_T^{1/T}\le\kappa\beta$.

For the lower bound, orient each graph edge once. A unit flow $F$ from
$o$ to $A_T$ has divergence $1$ at $o$, zero off the boundary, and total
divergence $-1$ on $A_T$. Its energy is
\[
 \mathcal E(F)=\sum_e\frac{F(e)^2}{c(e)}.
\]
Summation by parts and then Cauchy--Schwarz yield
\[
 1=\sum_e F(e)(v(\operatorname{tail}e)-v(\operatorname{head}e)),
 \qquad 1\le\mathcal E(F)\capT.
\]
Thus any such flow proves $\capT\ge1/\mathcal E(F)$.

\section{Concatenated disks and the lower bound}

We construct many growth paths with controlled congestion by appending
arbitrary disks in disjoint columns. Fix an integer $m\ge1$. Normalize
each translation class of $m$-square clumps so that the lowest square in
its leftmost occupied column is $Q_0$. Denote these representatives by
$D_1,\ldots,D_N$, where $N=a_m$, and put $L=m+1$.
For each $D_i$ choose one legal growth path from its port $Q_0$ to $D_i$,
retaining that port throughout.

\begin{lemma}\label{lem:join}
Let $(r,s)$ be the lowest center in the rightmost occupied column of a
clump $P$ containing $Q_0$. Then
\[
 P*D_i=P\cup\{Q_{(r+1,s)}\}\cup(D_i+(r+2,s))
\]
is a clump. It can be reached from $P$ by exactly $L$ legal additions,
all strictly to the right of the rightmost column of $P$.
\end{lemma}

\begin{proof}
Add the bridge $B=Q_{(r+1,s)}$, then the new port
$Q_{(r+2,s)}$, and then follow the chosen growth path of the translated
block. The bridge meets each component in one full vertical port edge;
possible contacts with an above-port square are endpoints of that same
edge. No square lies below either port in its extreme column, and the
old component and new block have no contact with each other.

At every stage the union is regular. Away from the bridge edges this is
the regularity of the individual components. At a lower bridge endpoint
the occupied sectors are two adjacent quadrants. At an upper endpoint
they are two adjacent quadrants or three consecutive quadrants, according
to whether the square above the port is present. Each punctured local
neighborhood is contractible; points inside a shared edge become interior.

The partial block retracts to its port by reversing its rooted growth
sequence. Each homotopy fixes the port, so it extends by the identity on
$P\cup\{B\}$. The port and then the bridge retract horizontally onto
their attached edges. Hence the union retracts onto $P$ and is contractible.

These are also legal single-square moves. For an internal block deletion,
extend its strong deformation retraction by the identity on the old
component and bridge. Any contact of the deleted square with the bridge
is in the retained port, so the definitions agree on the intersection.
Finite closed-set pasting gives a continuous homotopy, and the homotopy
still stays inside the deleted square. The bridge and initial-port moves
use the explicit horizontal retractions just described. Every intermediate
union and every deletion result is a clump. Since $P$ contains $Q_0$,
we have $r\ge0$, and every newly added square has first coordinate larger
than $r$. Thus none is $Q_0$; extending the translated block homotopies
by the identity on $P$ retains the actual pinned square. There are
$1+m=L$ additions.
\end{proof}

Starting from $o$, a completed word of $j$ blocks has size $n_j=1+jL$.
The following observation permits shared paths within a block while
preventing confusion between different completed parents.

\begin{lemma}\label{lem:prefix}
Different block words give different completed clumps. During stage $j$,
growth paths from different completed parents of size $n_j$ have no
common state or upward edge.
\end{lemma}

\begin{proof}
Every clump in these paths occupies consecutive columns beginning at zero.
For such a clump $C$, define its cumulative column count by
\[
 S_C(x)=\#\{Q_{(u,v)}\in C:u\le x\}.
\]
This count strictly increases across occupied integer columns. If $C$
lies on a stage-$j$ path after parent $P$, all squares added in that
stage have larger first coordinate than the rightmost column $r$ of $P$.
Thus $S_C(r)=n_j$, and this is the unique occupied column with that
count. Restricting $C$ to columns at most $r$ recovers its parent.
The lower endpoint of an upward edge therefore identifies its stage
and parent as well.

For a completed word, repeat this recovery at counts
$n_0,\ldots,n_{j-1}$. Between two successive cuts the first column is the
bridge, and the remaining $m$ squares give the normalized block after
translation. Thus the entire word is recovered uniquely.
\end{proof}

We now send a unit flow to area $n_k=1+kL$. Choose $k$ block labels
independently and uniformly from $\{1,\ldots,N\}$, and assign each
resulting monotone path mass $N^{-k}$. Summing these path flows gives a
unit flow $F_k$ in the size-$n_k$ stopped graph; set it equal to zero on
unused edges.

There are $N^j$ distinct stage-$j$ parents, each receiving mass $N^{-j}$.
Within one parent's stage, every edge has flow at most $N^{-j}$ and the
sum of edge flows is $L N^{-j}$, since every continuation has $L$ edges.
Therefore its sum of squared edge flows is at most $L N^{-2j}$.
Lemma~\ref{lem:prefix} makes the edge sets for different parents disjoint.
This estimate allows different next-block choices within the same parent
to share or merge their microscopic growth paths.

The lower size of a stage-$j$ edge lies between $1+jL$ and $(j+1)L$.
Since $\beta<1$, equation~\eqref{eq:conductance} gives
$c(e)\ge\beta^{(j+1)L}$. Consequently
\begin{equation}\label{eq:flow-energy}
 \mathcal E(F_k)
 \le L\beta^{-L}\sum_{j=0}^{k-1}(N\beta^L)^{-j}.
\end{equation}

\begin{proof}[Proof of Theorem~\ref{thm:main}]
Set $r_m=a_m\beta^{m+1}$. Its $(m+1)$st root tends to
$\kappa\beta<1$, so $r_m<1$ for every sufficiently large $m$.
Fix any such $m$, and write $r=r_m$. The geometric sum satisfies
\[
 \sum_{j=0}^{k-1}r^{-j}\le\frac{r^{-(k-1)}}{1-r}.
\]
The unit-flow bound, equations~\eqref{eq:escape} and
\eqref{eq:flow-energy}, give
\begin{equation}\label{eq:quantitative}
 p_{1+kL}\ge
 \frac{(1-r)\beta^{L-1}}{4L}\,r^{k-1},\qquad k\ge1.
\end{equation}
For any integer $T\ge2$, take $k=\lceil(T-1)/L\rceil$. Since
$1+kL\ge T$, the decreasing tail satisfies $p_T\ge p_{1+kL}$.
The positive prefactor in~\eqref{eq:quantitative} is fixed, and
$(k-1)/T\to1/L$. Taking $T$th roots yields
\[
 \liminf_{T\to\infty}p_T^{1/T}
 \ge r_m^{1/(m+1)}=\beta a_m^{1/(m+1)}.
\]
This holds for each sufficiently large fixed $m$. Letting $m\to\infty$
improves the same lower bound to $\kappa\beta$. Together with
\eqref{eq:upper}, it proves~\eqref{eq:rate} for the arbitrary fixed
admissible $\beta$. Finally, $\mu^{1/T}\to1$ because $0<\mu<\infty$,
so $\eta_T^{1/T}=\mu^{1/T}/p_T^{1/T}\to1/(\kappa\beta)$.
\end{proof}

\paragraph{AI assistance.}
OpenAI Codex assisted with argument development, checking cited sources,
and preparing the manuscript. The submitted proof is analytic and uses no
computational data. Independent mathematical review is requested.

\begin{thebibliography}{1}
\bibitem{BO}
Y.~Baryshnikov and E.~Onaran,
\emph{Ising disks: topology preserving Glauber dynamics},
Journal of Applied and Computational Topology \textbf{10} (2026), article~16.
\href{https://doi.org/10.1007/s41468-026-00248-w}{doi:10.1007/s41468-026-00248-w}.
Author preprint: \href{https://arxiv.org/abs/2410.22611}{arXiv:2410.22611}.
\end{thebibliography}
\end{document}
